\documentclass[12pt,a4paper]{article}
\usepackage[utf8]{inputenc}
\usepackage[margin=1in]{geometry}
\usepackage{graphicx}
\usepackage{enumitem}
\usepackage{hyperref}
\usepackage{xcolor}
\usepackage{longtable}
\usepackage{booktabs}
\usepackage{fancyhdr}
\usepackage{titlesec}

\definecolor{paleoblue}{RGB}{74,158,255}
\definecolor{paleogreen}{RGB}{76,175,80}
\definecolor{paleored}{RGB}{200,50,50}
\definecolor{darkgray}{RGB}{80,80,80}

\hypersetup{
  colorlinks=true,
  linkcolor=paleoblue,
  urlcolor=paleoblue,
}

\pagestyle{fancy}
\fancyhf{}
\fancyhead[L]{\small PaleoVox QA Test Document}
\fancyhead[R]{\small v1.0.6}
\fancyfoot[C]{\thepage}

\titleformat{\section}{\Large\bfseries}{\thesection}{1em}{}[\color{paleoblue}\titlerule]
\titleformat{\subsection}{\large\bfseries}{\thesubsection}{1em}{}

\newenvironment{testcase}[4]{%
  \smallskip
  \noindent\textbf{Test ID:} #1 \hfill \textbf{Area:} #2 \\
  \textbf{Title:} #3 \hfill \textbf{Priority:} #4
  \begin{description}[leftmargin=2em,style=nextline]
    \item[\textbf{Preconditions:}]
}{%
    \item[\textbf{Steps:}]
    \item[\textbf{Expected Result:}]
    \item[\textbf{Status:}] \rule{3cm}{0.4pt} \hspace{1em} \textbf{Date:} \rule{3cm}{0.4pt}
    \item[\textbf{Notes:}]
  \end{description}
  \smallskip
}

\begin{document}

\begin{titlepage}
\centering
\vspace*{3cm}

{\Huge\bfseries PaleoVox}\\[0.5cm]
{\Large Quality Assurance Test Document}\\[1.5cm]

{\large Fossil Voxel Augmentation GUI}\\[0.5cm]
{\large Version 1.0.6}\\[2cm]

\begin{tabular}{rl}
\textbf{Date:} & \today \\
\textbf{Authors:} & Alan Gabriel Amaro Colin \\
                 & Ángel Ángeles Cortés \\
                 & Jair Israel Barrientos Lara \\
                 & Jesus Alberto Díaz Cruz \\
\textbf{Institution:} & Facultad de Ciencias \& Instituto de Geología, UNAM \\
\end{tabular}

\vfill
\end{titlepage}

\tableofcontents
\newpage

% =============================================================================
\section{Introduction}
% =============================================================================

\subsection{Purpose}
This document defines the Quality Assurance test plan for the PaleoVox graphical user interface (GUI). Its purpose is to ensure that all user-facing functionality operates correctly, consistently, and handles errors gracefully before each release.

\subsection{Scope}
The test plan covers the entire PaleoVox GUI, including:
\begin{itemize}
  \item File loading via drag-and-drop, browse dialog, and the Load Mesh button
  \item Support for \texttt{.ply}, \texttt{.obj}, and \texttt{.npy} file formats
  \item Pipeline workflow: mesh loading $\rightarrow$ voxel conversion $\rightarrow$ morphological operations $\rightarrow$ mesh reconstruction
  \item Four augmentation operations: deformation, erosion, rotation, and fracture
  \item Visualization: 3D mesh view, 3D voxel view, mesh comparison, voxel comparison
  \item Analysis: t-SNE dimensionality reduction, 2D perspective projections
  \item File saving: meshes (\texttt{.ply}, \texttt{.obj}) and voxels (\texttt{.npy})
  \item Button state logic (enabled/disabled based on workflow state)
  \item Error handling and user feedback
  \item Reset and About functionality
\end{itemize}

\subsection{Test Environment}
\begin{itemize}
  \item \textbf{Operating System:} Linux (X11 session)
  \item \textbf{Python:} 3.8+ with PyQt5, Open3D, NumPy, SciPy, Matplotlib, scikit-learn
  \item \textbf{Window Manager:} Any X11-compatible WM (GNOME, KDE, XFCE, etc.)
  \item \textbf{Screen Resolution:} 1100$\times$700 minimum
  \item \textbf{Test Data:} Representative \texttt{.ply}, \texttt{.obj}, and \texttt{.npy} files
\end{itemize}

\newpage

% =============================================================================
\section{Test Case Catalog}
% =============================================================================

% -----------------------------------------------------------------------------
\subsection{File Loading \& Drop Zone}
% -----------------------------------------------------------------------------

\begin{testcase}{FL-001}{File Loading}{Drag \& drop \texttt{.ply} file onto Drop Zone}{High}
  \item No file is currently loaded.
  \item Drag a valid \texttt{.ply} file from the file manager onto the Drop Zone.
  \item The Drop Zone border turns blue on hover, then green on drop. The file loads; File Info panel shows vertex/triangle counts. Status bar confirms load.
\end{testcase}

\begin{testcase}{FL-002}{File Loading}{Drag \& drop \texttt{.obj} file onto Drop Zone}{High}
  \item No file is currently loaded.
  \item Drag a valid \texttt{.obj} file from the file manager onto the Drop Zone.
  \item The Drop Zone accepts the file; status turns green. Mesh loads correctly. File Info updated.
\end{testcase}

\begin{testcase}{FL-003}{File Loading}{Drag \& drop \texttt{.npy} file onto Drop Zone}{High}
  \item No file is currently loaded.
  \item Drag a valid \texttt{.npy} file (3D voxel array) onto the Drop Zone.
  \item Drop Zone accepts the file. Voxel array loads (no mesh). File Info shows voxel shape and occupied count. Mesh vertex/triangle labels show ``---". Augmentation buttons enabled; mesh-dependent buttons disabled.
\end{testcase}

\begin{testcase}{FL-004}{File Loading}{Drag \& drop unsupported file type onto Drop Zone}{Medium}
  \item No file is currently loaded.
  \item Drag a \texttt{.txt}, \texttt{.png}, or other unsupported file onto the Drop Zone.
  \item The file is rejected (ignored). Drop Zone remains in idle state. No error or status change.
\end{testcase}

\begin{testcase}{FL-005}{File Loading}{Browse button opens file dialog with correct filters}{High}
  \item No file is currently loaded.
  \item Click the ``Browse...'' button.
  \item File dialog opens with title ``Open 3D Mesh or Voxel Grid''. Filter options include ``Supported Files (*.ply *.obj *.npy)'', ``3D Mesh Files'', ``PLY'', ``OBJ'', and ``Voxel Grid Files''.
\end{testcase}

\begin{testcase}{FL-006}{File Loading}{Load via Browse selecting a \texttt{.ply} file}{High}
  \item No file is currently loaded.
  \item Click Browse, select a valid \texttt{.ply} file, click Open.
  \item File loads successfully. Same behavior as FL-001.
\end{testcase}

\begin{testcase}{FL-007}{File Loading}{Load via Browse selecting a \texttt{.npy} file}{High}
  \item No file is currently loaded.
  \item Click Browse, select a valid \texttt{.npy} voxel file, click Open.
  \item Voxel array loads. Same behavior as FL-003.
\end{testcase}

\begin{testcase}{FL-008}{File Loading}{Click on Drop Zone text opens file dialog}{Medium}
  \item No file is currently loaded.
  \item Click directly on the Drop Zone text/label area.
  \item File dialog opens (same as Browse button behavior).
\end{testcase}

\begin{testcase}{FL-009}{File Loading}{Load Mesh button re-loads current file}{Low}
  \item A mesh file is already loaded (e.g., via Browse).
  \item Go to Pipeline tab, click ``Load Mesh''.
  \item The same file is reloaded. File Info and button states are refreshed.
\end{testcase}

\begin{testcase}{FL-010}{File Loading}{Drag \& drop visual state transitions}{Medium}
  \item No file is currently loaded.
  \item Hover over the Drop Zone with a valid \texttt{.ply} file (do not drop); then drag away; then drop the file.
  \item Hover $\rightarrow$ blue dashed border. Drag leave $\rightarrow$ gray idle border. Drop $\rightarrow$ green dashed border with file name displayed.
\end{testcase}

\begin{testcase}{FL-011}{File Loading}{Load corrupted/invalid file}{High}
  \item A test file that is not a valid mesh or voxel grid is available.
  \item Attempt to load the corrupted file via Browse.
  \item An error dialog appears with title ``Load Error'' and a descriptive error message. Application does not crash. Previous state preserved.
\end{testcase}

% -----------------------------------------------------------------------------
\subsection{File Info Panel}
% -----------------------------------------------------------------------------

\begin{testcase}{FI-001}{File Info}{Info panel shows mesh statistics after \texttt{.ply} load}{Medium}
  \item Load a valid \texttt{.ply} mesh file.
  \item Observe the File Info panel on the left sidebar.
  \item ``Path'' shows the file path (home abbreviated as \texttt{\~{}}). ``Vertices'' and ``Triangles'' show correct counts. ``Voxel shape'' and ``Occupied voxels'' show ``---''.
\end{testcase}

\begin{testcase}{FI-002}{File Info}{Info panel shows voxel statistics after \texttt{.npy} load}{Medium}
  \item Load a valid \texttt{.npy} voxel file (e.g., $(128,128,128)$ with $15243$ occupied).
  \item Observe the File Info panel.
  \item ``Vertices'' and ``Triangles'' show ``---''. ``Voxel shape'' shows the shape tuple. ``Occupied voxels'' shows correct count.
\end{testcase}

\begin{testcase}{FI-003}{File Info}{Info panel updated after voxelization}{Medium}
  \item Load a mesh, then click ``Mesh $\rightarrow$ Voxels''.
  \item Observe the File Info panel.
  \item ``Voxel shape'' and ``Occupied voxels'' are populated. Vertices/triangles remain from the mesh.
\end{testcase}

% -----------------------------------------------------------------------------
\subsection{Pipeline Tab}
% -----------------------------------------------------------------------------

\begin{testcase}{PL-001}{Pipeline}{Mesh $\rightarrow$ Voxels with default parameters}{High}
  \item A mesh file is loaded.
  \item Set npoints=100000, dim=256. Click ``Mesh $\rightarrow$ Voxels''.
  \item Voxel grid of shape $(256,256,256)$ is created. Status bar confirms conversion. Augmentation buttons become enabled. Original voxel saved for comparison.
\end{testcase}

\begin{testcase}{PL-002}{Pipeline}{Mesh $\rightarrow$ Voxels with custom parameters}{Medium}
  \item A mesh file is loaded.
  \item Set npoints=50000, dim=64. Click ``Mesh $\rightarrow$ Voxels''.
  \item Voxel grid of shape $(64,64,64)$ is created. Occupied voxel count varies with parameters.
\end{testcase}

\begin{testcase}{PL-003}{Pipeline}{Mesh $\rightarrow$ Voxels boundary values}{Low}
  \item A mesh file is loaded.
  \item Test npoints=100 (min), npoints=200000 (max), dim=32 (min), dim=512 (max).
  \item All boundary values produce a valid voxel grid without crashes or hangs.
\end{testcase}

\begin{testcase}{PL-004}{Pipeline}{Dilate Voxels with default iterations}{High}
  \item Voxels exist (after Mesh $\rightarrow$ Voxels).
  \item Set iterations=2. Click ``Dilate Voxels''.
  \item Voxel array is morphologically closed. Occupied voxel count increases. Status bar confirms. Mesh is NOT updated until explicit reconstruction.
\end{testcase}

\begin{testcase}{PL-005}{Pipeline}{Dilate Voxels with various iterations}{Medium}
  \item Voxels exist.
  \item Test iterations=1, 5, and 10.
  \item Higher iterations produce more aggressive filling. No error at boundary values.
\end{testcase}

\begin{testcase}{PL-006}{Pipeline}{Voxels $\rightarrow$ Mesh reconstruction}{High}
  \item Voxels exist (after conversion or \texttt{.npy} load).
  \item Click ``Voxels $\rightarrow$ Mesh''.
  \item A reconstructed triangle mesh is generated. Vertices and triangles displayed in File Info and Reconstruction info label. Status bar confirms.
\end{testcase}

\begin{testcase}{PL-007}{Pipeline}{Full pipeline workflow (end-to-end)}{High}
  \item No file loaded.
  \item Load \texttt{.ply} mesh $\rightarrow$ Mesh $\rightarrow$ Voxels $\rightarrow$ Dilate $\rightarrow$ Voxels $\rightarrow$ Mesh.
  \item All steps complete without errors. Final mesh is reconstructed from dilated voxels. File Info updated at each step.
\end{testcase}

\begin{testcase}{PL-008}{Pipeline}{Pipeline buttons disabled when no mesh/voxels}{High}
  \item No file loaded (fresh start).
  \item Observe Pipeline tab buttons.
  \item All buttons except ``Load Mesh'' are disabled. After loading mesh: Load Mesh, Mesh $\rightarrow$ Voxels enabled; Dilate, Voxels $\rightarrow$ Mesh disabled. After voxelization: all enabled.
\end{testcase}

% -----------------------------------------------------------------------------
\subsection{Augmentations}
% -----------------------------------------------------------------------------

\begin{testcase}{AU-001}{Augmentation}{Deformation on each axis}{High}
  \item Voxels exist. original\_voxel has been captured.
  \item For each axis (X, Y, Z), set factor=0.85. Click Apply next to Deformation.
  \item Occupied voxel count changes. ``Current'' count in Voxel Comparison updates. Original voxel count unchanged. Status bar confirms.
\end{testcase}

\begin{testcase}{AU-002}{Augmentation}{Deformation with min/max factor}{Medium}
  \item Voxels exist.
  \item Test factor=0.10 (min, heavy compression) and factor=0.99 (max, minimal compression).
  \item Both produce valid results without error. Heavier compression reduces occupied count significantly.
\end{testcase}

\begin{testcase}{AU-003}{Augmentation}{Erosion on each axis}{High}
  \item Voxels exist.
  \item For X, Y, Z axes: set increment=0.5. Click Apply next to Erosion.
  \item Occupied voxel count decreases on each axis. Random cut removes a portion. Current voxel count updates.
\end{testcase}

\begin{testcase}{AU-004}{Augmentation}{Erosion with min/max increment}{Medium}
  \item Voxels exist.
  \item Test increment=0.01 (min, almost no removal) and increment=1.0 (max, heavy removal).
  \item Both produce valid results. Min increment preserves most voxels. Max removes a large portion.
\end{testcase}

\begin{testcase}{AU-005}{Augmentation}{Rotation with zero angles}{Low}
  \item Voxels exist.
  \item Set all rotation angles (X, Y, Z) to 0. Click Apply.
  \item Voxel array is unchanged (identity rotation). Occupied count should match.
\end{testcase}

\begin{testcase}{AU-006}{Augmentation}{Rotation with non-zero angles}{Medium}
  \item Voxels exist.
  \item Set X=30°, Y=45°, Z=90°. Click Apply.
  \item Rotation is applied. Occupied count may change slightly due to nearest-neighbor interpolation. No errors.
\end{testcase}

\begin{testcase}{AU-007}{Augmentation}{Rotation with max angles}{Low}
  \item Voxels exist.
  \item Set all angles to 180°. Click Apply.
  \item Rotation completes. Object should be flipped (upside-down).
\end{testcase}

\begin{testcase}{AU-008}{Augmentation}{Fracture with return\_both=False}{High}
  \item Voxels exist.
  \item Uncheck ``Return both''. Set max\_pos=10. Click Apply.
  \item Fracture pattern is subtracted from voxel array. Occupied count decreases. Status bar shows count only.
\end{testcase}

\begin{testcase}{AU-009}{Augmentation}{Fracture with return\_both=True}{Medium}
  \item Voxels exist.
  \item Check ``Return both''. Set max\_pos=5. Click Apply.
  \item voxel and fracture\_pattern are separated. Status bar shows pattern voxel count. Fracture lines are isolated.
\end{testcase}

\begin{testcase}{AU-010}{Augmentation}{Fracture max\_pos boundary values}{Low}
  \item Voxels exist.
  \item Test max\_pos=1 (straight fracture) and max\_pos=50 (random wandering).
  \item Both produce valid results. Low value = more directed; high value = more random.
\end{testcase}

\begin{testcase}{AU-011}{Augmentation}{Save Deformed Voxels to \texttt{.npy}}{High}
  \item Voxels exist (preferably after an augmentation).
  \item Click ``Save Deformed Voxels''. Choose a file path and save.
  \item Voxel array is saved as \texttt{.npy}. Status bar confirms the path. File can be reloaded successfully.
\end{testcase}

\begin{testcase}{AU-012}{Augmentation}{Augmentations chain sequentially}{Medium}
  \item Voxels exist.
  \item Apply Deformation $\rightarrow$ Erosion $\rightarrow$ Rotation $\rightarrow$ Fracture.
  \item Each augmentation operates on the current voxel state. All complete without error. Final state reflects all augmentations.
\end{testcase}

\begin{testcase}{AU-013}{Augmentation}{Augmentation buttons disabled without voxels}{Medium}
  \item Load a mesh only (no voxelization).
  \item Observe augmentation controls.
  \item All Apply buttons and ``Save Deformed Voxels'' are disabled.
\end{testcase}

% -----------------------------------------------------------------------------
\subsection{View \& Save Tab}
% -----------------------------------------------------------------------------

\begin{testcase}{VS-001}{View/Save}{View Mesh opens Open3D viewer}{High}
  \item A mesh is loaded (from \texttt{.ply} or after reconstruction).
  \item Select a color from the dropdown (e.g., Red). Click ``View Mesh''.
  \item An Open3D viewer window opens (1024$\times$768) showing the mesh in the chosen color. GUI remains responsive.
\end{testcase}

\begin{testcase}{VS-002}{View/Save}{View Voxels opens Open3D viewer}{High}
  \item Voxels exist.
  \item Select Cyan color. Click ``View Voxels''.
  \item An Open3D viewer shows the voxel grid as a point cloud in the chosen color. GUI remains responsive.
\end{testcase}

\begin{testcase}{VS-003}{View/Save}{View Voxels with no occupied voxels}{Low}
  \item An empty voxel array exists (all zeros). (May require corner-case setup.)
  \item Click ``View Voxels''.
  \item Status bar shows ``No occupied voxels to display''. No viewer opens. No error dialog.
\end{testcase}

\begin{testcase}{VS-004}{View/Save}{Color picker changes display color}{Medium}
  \item Voxels or mesh exist.
  \item Cycle through each color in the dropdown. View Mesh or View Voxels after each change.
  \item Each color is applied correctly in the viewer. Color selection persists for subsequent views.
\end{testcase}

\begin{testcase}{VS-005}{View/Save}{Save Mesh to \texttt{.ply}}{High}
  \item A mesh exists.
  \item Click ``Save Mesh''. Choose a path ending in \texttt{.ply}. Save.
  \item Mesh is saved. Status bar confirms. File can be reopened.
\end{testcase}

\begin{testcase}{VS-006}{View/Save}{Save Mesh to \texttt{.obj}}{Medium}
  \item A mesh exists.
  \item Click ``Save Mesh''. Choose a path ending in \texttt{.obj}. Save.
  \item Mesh is saved as \texttt{.obj}. File is valid.
\end{testcase}

\begin{testcase}{VS-007}{View/Save}{Save Voxel to \texttt{.npy}}{High}
  \item Voxels exist.
  \item Click ``Save Voxel''. Choose a path. Save.
  \item Voxel array saved as \texttt{.npy}. Status bar confirms. Can be reloaded with FL-003/FL-007.
\end{testcase}

\begin{testcase}{VS-008}{View/Save}{Save with default filename derived from source}{Low}
  \item A mesh file named \texttt{fossil.ply} is loaded.
  \item Click ``Save Mesh'' and observe the suggested filename, then ``Save Voxel'' (after voxelization).
  \item Save Mesh suggests \texttt{fossil\_processed.ply}. Save Voxel suggests \texttt{fossil\_voxel.npy}.
\end{testcase}

% -----------------------------------------------------------------------------
\subsection{Reconstruction \& Comparison Tab}
% -----------------------------------------------------------------------------

\begin{testcase}{RC-001}{Recon/Comp}{Reconstruct from Voxels}{High}
  \item Voxels exist but no reconstructed mesh.
  \item Click ``Reconstruct from Voxels''.
  \item Mesh is reconstructed. Reconstruction info label updated (vertices, triangles). ``Save Reconstructed Mesh'' button enabled. ``Compare Original vs Reconstructed'' enabled if original mesh exists.
\end{testcase}

\begin{testcase}{RC-002}{Recon/Comp}{Reconstruct from \texttt{.npy}-loaded voxels}{High}
  \item A \texttt{.npy} file is loaded (no original mesh).
  \item Click ``Reconstruct from Voxels''.
  \item Reconstruction succeeds. Since no original mesh, ``Compare Original vs Reconstructed'' remains disabled. Mesh view and save work.
\end{testcase}

\begin{testcase}{RC-003}{Recon/Comp}{Compare Original vs Reconstructed --- Both}{High}
  \item Original mesh and reconstructed mesh both exist.
  \item Select ``Both'' from the show dropdown. Click ``Compare Original vs Reconstructed''.
  \item Open3D viewer opens with original (blue) and reconstructed (red) overlaid. Title shows ``Original (Blue) vs Reconstructed (Red)''.
\end{testcase}

\begin{testcase}{RC-004}{Recon/Comp}{Compare Original vs Reconstructed --- Original Only}{Medium}
  \item Original mesh and reconstructed mesh both exist.
  \item Select ``Original Only''. Click compare.
  \item Viewer shows only the original mesh in blue. Title reflects ``Original Only''.
\end{testcase}

\begin{testcase}{RC-005}{Recon/Comp}{Compare Original vs Reconstructed --- Reconstructed Only}{Medium}
  \item Original mesh and reconstructed mesh both exist.
  \item Select ``Reconstructed Only''. Click compare.
  \item Viewer shows only the reconstructed mesh in red.
\end{testcase}

\begin{testcase}{RC-006}{Recon/Comp}{Compare Voxels --- Both}{High}
  \item Original voxels and current voxels both exist (after an augmentation).
  \item Select ``Both''. Click ``Compare Voxels''.
  \item Open3D viewer shows original (blue) and current (red) voxel grids overlaid. Voxel counts in info label are correct.
\end{testcase}

\begin{testcase}{RC-007}{Recon/Comp}{Compare Voxels --- Original Only}{Medium}
  \item Original voxels and current voxels both exist.
  \item Select ``Original Only''. Click compare.
  \item Only original voxels shown in blue.
\end{testcase}

\begin{testcase}{RC-008}{Recon/Comp}{Compare Voxels --- Current Only}{Medium}
  \item Original voxels and current voxels both exist.
  \item Select ``Current Only''. Click compare.
  \item Only current voxels shown in red.
\end{testcase}

\begin{testcase}{RC-009}{Recon/Comp}{Voxel comparison info label}{Medium}
  \item Original voxels exist (10000 occupied). Apply erosion. Current voxels now have 8000 occupied.
  \item Observe the ``Original: --- Current: ---'' label in Voxel Comparison group.
  \item Label reads ``Original: 10,000 voxels   Current: 8,000 voxels''.
\end{testcase}

\begin{testcase}{RC-010}{Recon/Comp}{Save Reconstructed Mesh}{High}
  \item A reconstructed mesh exists.
  \item Click ``Save Reconstructed Mesh''. Choose a path and save.
  \item Mesh saved. Status bar confirms. Default filename derived from source.
\end{testcase}

% -----------------------------------------------------------------------------
\subsection{t-SNE Analysis Tab}
% -----------------------------------------------------------------------------

\begin{testcase}{TS-001}{t-SNE}{Generate t-SNE with default parameters}{High}
  \item Both original\_voxel and voxel exist and have occupied voxels.
  \item Use default params (perplexity=100, seed=42, percentage=0.50, size=1.0). Click ``Generate t-SNE''.
  \item A dialog opens showing the t-SNE scatter plot (two side-by-side subplots). Dialog has Save Image and Close buttons. Status bar confirms completion.
\end{testcase}

\begin{testcase}{TS-002}{t-SNE}{Generate t-SNE with custom parameters}{Medium}
  \item Voxels exist.
  \item Set perplexity=30, seed=9999, percentage=0.75, size=2.0. Click generate.
  \item t-SNE generates with custom parameters. Different perplexity produces different embedding. No error.
\end{testcase}

\begin{testcase}{TS-003}{t-SNE}{t-SNE parameter boundary values}{Low}
  \item Voxels exist.
  \item Test perplexity=5 (min), perplexity=500 (max), percentage=0.01 (min), percentage=1.00 (max).
  \item All boundary values produce valid results without crash.
\end{testcase}

\begin{testcase}{TS-004}{t-SNE}{Error when no occupied voxels}{Medium}
  \item Voxels exist but both original and current have zero occupied voxels (edge case).
  \item Click ``Generate t-SNE''.
  \item Error dialog: ``Both original and current voxels must contain occupied voxels.'' No crash.
\end{testcase}

\begin{testcase}{TS-005}{t-SNE}{Save t-SNE image from dialog}{Medium}
  \item t-SNE dialog is open.
  \item Click ``Save Image...''. Choose a path (PNG). Save.
  \item Image saved to chosen path. Status bar confirms.
\end{testcase}

\begin{testcase}{TS-006}{t-SNE}{Close t-SNE dialog without saving}{Low}
  \item t-SNE dialog is open.
  \item Click ``Close''.
  \item Dialog closes. Temp file is cleaned up. GUI state unchanged.
\end{testcase}

\begin{testcase}{TS-007}{t-SNE}{Button disabled without original voxels}{Medium}
  \item Load a mesh and convert to voxels (original = current). Then load a new mesh.
  \item Observe ``Generate t-SNE'' button state after each state change.
  \item Button enabled only when both original\_voxel and voxel are not None.
\end{testcase}

% -----------------------------------------------------------------------------
\subsection{2D Perspectives Tab}
% -----------------------------------------------------------------------------

\begin{testcase}{D2-001}{2D Persp}{Generate single 2D perspective --- XY view}{High}
  \item Voxels exist.
  \item Select Axis=XY, Color=Blue, Marker=o, Size=1.0. Click ``Generate 2D Perspective''.
  \item Dialog opens showing the XY projection scatter plot. Axes labeled X and Y. Save Image and Close buttons present.
\end{testcase}

\begin{testcase}{D2-002}{2D Persp}{Generate single 2D perspective --- XZ view}{Medium}
  \item Voxels exist.
  \item Select Axis=XZ. Click generate.
  \item Dialog shows XZ projection. Axes labeled X and Z.
\end{testcase}

\begin{testcase}{D2-003}{2D Persp}{Generate single 2D perspective --- YZ view}{Medium}
  \item Voxels exist.
  \item Select Axis=YZ. Click generate.
  \item Dialog shows YZ projection. Axes labeled Y and Z.
\end{testcase}

\begin{testcase}{D2-004}{2D Persp}{Single perspective with different colors}{Low}
  \item Voxels exist.
  \item Test each color option (Blue, Red, Green, Orange, Purple, Cyan, Yellow, White, Gray).
  \item Each color renders correctly in the scatter plot.
\end{testcase}

\begin{testcase}{D2-005}{2D Persp}{Single perspective with different markers}{Low}
  \item Voxels exist.
  \item Test each marker option (o, ., \^{}, s, *, +, x).
  \item Each marker renders correctly.
\end{testcase}

\begin{testcase}{D2-006}{2D Persp}{Single perspective marker size}{Low}
  \item Voxels exist.
  \item Test size=0.1 (min) and size=10.0 (max).
  \item Min size produces tiny points; max size produces large points. Both functional.
\end{testcase}

\begin{testcase}{D2-007}{2D Persp}{Generate comparison (original vs current)}{High}
  \item Both original\_voxel and voxel exist.
  \item Set Axis=XY. Customize colors/markers/sizes for both. Click ``Generate Comparison''.
  \item Dialog opens showing overlaid scatter plot. Original in color1/marker1, current in color2/marker2. Legend shows ``Original'' and ``Current''.
\end{testcase}

\begin{testcase}{D2-008}{2D Persp}{Comparison with different axes}{Medium}
  \item Voxels exist.
  \item Generate comparison for XZ and YZ views.
  \item Both projections work correctly. Axes are properly labeled.
\end{testcase}

\begin{testcase}{D2-009}{2D Persp}{Save 2D image from dialog}{Medium}
  \item 2D perspective dialog is open.
  \item Click ``Save Image...''. Choose path. Save.
  \item Image saved correctly. Status bar confirms.
\end{testcase}

\begin{testcase}{D2-010}{2D Persp}{Button states for 2D perspectives}{Medium}
  \item Test with: (a) no voxels, (b) voxels exist but no original, (c) both exist.
  \item (a) Both buttons disabled. (b) Single enabled, Comparison disabled. (c) Both enabled.
\end{testcase}

% -----------------------------------------------------------------------------
\subsection{Reset \& About}
% -----------------------------------------------------------------------------

\begin{testcase}{RA-001}{Reset/About}{Reset All clears all state}{High}
  \item A full workflow has been performed (mesh loaded, voxelized, augmented, reconstructed).
  \item Click ``Reset All''.
  \item All internal state cleared. All labels show ``---''. Drop Zone returns to idle style. All buttons disabled except Browse. Status bar shows ``Reset --- ready for new file''.
\end{testcase}

\begin{testcase}{RA-002}{Reset/About}{Reset All from \texttt{.npy} only state}{Medium}
  \item A \texttt{.npy} file is loaded. No mesh.
  \item Click ``Reset All''.
  \item All state cleared. Same as RA-001.
\end{testcase}

\begin{testcase}{RA-003}{Reset/About}{About dialog displays correctly}{Low}
  \item Application is running.
  \item Click ``More about PaleoVox...''.
  \item A dialog opens with PaleoVox description, creators list with names and emails. Close button works. Dialog is modal.
\end{testcase}

\begin{testcase}{RA-004}{Reset/About}{About dialog close button}{Low}
  \item About dialog is open.
  \item Click ``Close''.
  \item Dialog closes. Main window state unchanged.
\end{testcase}

% -----------------------------------------------------------------------------
\subsection{Status Bar \& Error Handling}
% -----------------------------------------------------------------------------

\begin{testcase}{SE-001}{Status/Error}{Status bar shows initial ready message}{Low}
  \item Fresh application start.
  \item Observe status bar at bottom.
  \item Shows ``Ready --- Drag \& drop a .ply, .obj or .npy file to begin''.
\end{testcase}

\begin{testcase}{SE-002}{Status/Error}{Status bar updates on each operation}{Medium}
  \item Perform various operations (load, voxelize, augment, reconstruct).
  \item Observe status bar after each.
  \item Each operation updates the status bar with a descriptive message including counts.
\end{testcase}

\begin{testcase}{SE-003}{Status/Error}{Error dialog on backend failure}{Medium}
  \item Induce an error (e.g., attempt to save to read-only directory).
  \item Trigger the operation.
  \item An error dialog appears with relevant title and message. Application does not crash. Status unchanged for unaffected state.
\end{testcase}

\begin{testcase}{SE-004}{Status/Error}{Window title and icon}{Low}
  \item Application is running.
  \item Observe the window title bar.
  \item Title: ``PaleoVox --- Fossil Voxel Augmentation''. Icon is the logo if available.
\end{testcase}

% =============================================================================
\newpage
\section{Integration Test Scenarios}
% =============================================================================

\subsection{End-to-End Workflow: Mesh File}
\begin{enumerate}[leftmargin=1.5em]
  \item Drag-and-drop a \texttt{.ply} file onto the Drop Zone. \hfill\textit{Expect: Loads with green style.}
  \item Verify File Info panel shows correct vertices/triangles. \hfill\textit{Expect: Populated.}
  \item Go to Pipeline tab. Set npoints=50000, dim=128. Click Mesh $\rightarrow$ Voxels. \hfill\textit{Expect: Voxel grid created.}
  \item Click Dilate Voxels (iterations=2). \hfill\textit{Expect: Occupied count increases.}
  \item Go to Augmentations. Apply Deformation (Z, factor=0.85). \hfill\textit{Expect: Voxels compacted.}
  \item Apply Erosion (Y, increment=0.3). \hfill\textit{Expect: Voxels partially removed.}
  \item Apply Rotation (X=15°, Y=0°, Z=0°). \hfill\textit{Expect: Voxels rotated.}
  \item Click Voxels $\rightarrow$ Mesh. \hfill\textit{Expect: Mesh reconstructed.}
  \item Go to Reconstruction tab. Click Compare Original vs Reconstructed (Both). \hfill\textit{Expect: Viewer shows overlay.}
  \item Go to View \& Save tab. Click Save Mesh. \hfill\textit{Expect: Saved successfully.}
  \item Click Save Voxel. \hfill\textit{Expect: Saved successfully.}
  \item Go to 2D Perspectives tab. Generate XY comparison. \hfill\textit{Expect: Dialog with overlay scatter.}
  \item Go to t-SNE Analysis tab. Generate t-SNE. \hfill\textit{Expect: Dialog with embedding plot.}
  \item Save t-SNE image. \hfill\textit{Expect: Save dialog and successful save.}
  \item Click Reset All. \hfill\textit{Expect: All cleared.}
\end{enumerate}

\subsection{End-to-End Workflow: Voxel (.npy) File}
\begin{enumerate}[leftmargin=1.5em]
  \item Click Browse. Select a \texttt{.npy} file. \hfill\textit{Expect: Loads with voxel info.}
  \item Verify File Info: shape and occupied voxels shown, vertices/triangles = ``---''. \hfill\textit{Expect: Correct.}
  \item Verify augmentation buttons are enabled (Deform, Erode, Rotate, Fracture). \hfill\textit{Expect: All enabled.}
  \item Verify mesh-dependent buttons (View Mesh, Save Mesh) are disabled. \hfill\textit{Expect: Disabled.}
  \item Apply Erosion (X, increment=0.5). \hfill\textit{Expect: Voxel count decreases.}
  \item Click Reconstruct from Voxels. \hfill\textit{Expect: Mesh reconstructed. View Mesh now enabled.}
  \item View the reconstructed mesh. \hfill\textit{Expect: Open3D viewer opens.}
  \item Save the reconstructed mesh as \texttt{.ply}. \hfill\textit{Expect: Saved.}
  \item Generate 2D single perspective (XZ view). \hfill\textit{Expect: Dialog with scatter plot.}
  \item Generate t-SNE comparison. \hfill\textit{Expect: Dialog with embedding.}
\end{enumerate}

\subsection{Sequential Augmentation Pipeline}
\begin{enumerate}[leftmargin=1.5em]
  \item Load mesh. Convert to voxels (dim=128).
  \item Record original voxel count. \hfill\textit{Note value.}
  \item Deformation (Z, 0.85) $\rightarrow$ Erosion (Y, 0.5) $\rightarrow$ Rotation (X=30°, Y=45°, Z=90°) $\rightarrow$ Fracture (max\_pos=10).
  \item Voxel count should differ from original after each step. \hfill\textit{Expect: Cumulative changes.}
  \item Compare original vs current voxels. \hfill\textit{Expect: Visible differences in blue/red overlay.}
  \item Compare t-SNE. \hfill\textit{Expect: Embeddings show distribution differences.}
\end{enumerate}

% =============================================================================
\section{Button State Matrix}
% =============================================================================

The following table summarizes expected button enabled states for key workflow stages.

\small
\begin{longtable}{lcccccc}
\toprule
\textbf{Button} & \textbf{Start} & \textbf{Mesh Only} & \textbf{Voxels} & \textbf{Reconstructed} & \textbf{npy Load} \\
\midrule
Load Mesh & $\square$ & $\blacksquare$ & $\blacksquare$ & $\blacksquare$ & $\blacksquare$ \\
Mesh $\rightarrow$ Voxels & $\square$ & $\blacksquare$ & $\square$ & $\square$ & $\square$ \\
Dilate Voxels & $\square$ & $\square$ & $\blacksquare$ & $\blacksquare$ & $\blacksquare$ \\
Voxels $\rightarrow$ Mesh & $\square$ & $\square$ & $\blacksquare$ & $\blacksquare$ & $\blacksquare$ \\
Deformation Apply & $\square$ & $\square$ & $\blacksquare$ & $\blacksquare$ & $\blacksquare$ \\
Erosion Apply & $\square$ & $\square$ & $\blacksquare$ & $\blacksquare$ & $\blacksquare$ \\
Rotation Apply & $\square$ & $\square$ & $\blacksquare$ & $\blacksquare$ & $\blacksquare$ \\
Fracture Apply & $\square$ & $\square$ & $\blacksquare$ & $\blacksquare$ & $\blacksquare$ \\
Save Deformed Voxels & $\square$ & $\square$ & $\blacksquare$ & $\blacksquare$ & $\blacksquare$ \\
View Mesh & $\square$ & $\blacksquare$ & $\square$ & $\blacksquare$ & $\square$ \\
View Voxels & $\square$ & $\square$ & $\blacksquare$ & $\blacksquare$ & $\blacksquare$ \\
Save Mesh & $\square$ & $\blacksquare$ & $\square$ & $\blacksquare$ & $\square$ \\
Save Voxel & $\square$ & $\square$ & $\blacksquare$ & $\blacksquare$ & $\blacksquare$ \\
Reconstruct from Voxels & $\square$ & $\square$ & $\blacksquare$ & $\blacksquare$ & $\blacksquare$ \\
Compare Orig vs Recon & $\square$ & $\square$ & $\square$ & $\blacksquare$* & $\square$ \\
Save Reconstructed & $\square$ & $\square$ & $\square$ & $\blacksquare$ & $\square$ \\
Compare Voxels & $\square$ & $\square$ & $\blacksquare$** & $\blacksquare$** & $\blacksquare$** \\
Generate t-SNE & $\square$ & $\square$ & $\blacksquare$** & $\blacksquare$** & $\blacksquare$** \\
2D Single Perspective & $\square$ & $\square$ & $\blacksquare$ & $\blacksquare$ & $\blacksquare$ \\
2D Comparison & $\square$ & $\square$ & $\blacksquare$** & $\blacksquare$** & $\blacksquare$** \\
\bottomrule
\multicolumn{7}{l}{\footnotesize $\blacksquare$ = enabled; $\square$ = disabled} \\
\multicolumn{7}{l}{\footnotesize * Requires original\_mesh AND reconstructed\_mesh} \\
\multicolumn{7}{l}{\footnotesize ** Requires original\_voxel AND voxel (both not None)} \\
\multicolumn{7}{l}{\footnotesize Start = no file loaded; Mesh Only = .ply/.obj loaded; Voxels = voxelized; Reconstructed = mesh rebuilt; npy Load = .npy file loaded}
\end{longtable}
\normalsize

% =============================================================================
\section{Error Handling Coverage}
% =============================================================================

\begin{longtable}{p{4cm}p{8cm}p{3cm}}
\toprule
\textbf{Error Scenario} & \textbf{Expected Behavior} & \textbf{Test ID} \\
\midrule
Load corrupted/invalid file & Error dialog. No crash. Previous state preserved. & FL-011 \\
Load unsupported file type & File rejected silently (drag/drop) or not shown (dialog filter). & FL-004 \\
Mesh $\rightarrow$ Voxels with no mesh & Handler returns early. No error. & --- \\
Augmentation with no voxels & Handler returns early. No error. & AU-013 \\
View empty voxel grid & Status bar message. No viewer opened. & VS-003 \\
t-SNE with zero occupied voxels & Error dialog. No crash. & TS-004 \\
Reconstruct from empty voxels & Backend handles. Error reported gracefully. & --- \\
Save to read-only directory & Error dialog with details. & SE-003 \\
Load non-existent file path & File dialog prevents selection. & --- \\
Threaded operations (viewers) & GUI remains responsive during 3D viewer. & VS-001, VS-002 \\
Temp file cleanup after dialog close & No orphaned temp files accumulate. & TS-006, D2-009 \\
\bottomrule
\end{longtable}

% =============================================================================
\section{Regression Tests}
% =============================================================================

\begin{enumerate}[leftmargin=1.5em]
  \item \textbf{RT-001}: After \texttt{.npy} loading was added, verify that \texttt{.ply} and \texttt{.obj} loading still works identically. \hfill\textit{(FL-001, FL-002)}
  \item \textbf{RT-002}: After 2D Perspectives tab was added, verify that existing tabs (Pipeline, Augment, Reconstruction, View \& Save, t-SNE) still function correctly. \hfill\textit{(All existing tests)}
  \item \textbf{RT-003}: After \texttt{save\_path} parameter was added to \texttt{plot\_2d\_perspective}, verify that standalone backend function still works (calling without \texttt{save\_path} opens \texttt{plt.show()}). \hfill\textit{(Manual backend test)}
  \item \textbf{RT-004}: Drop Zone style transitions (idle $\rightarrow$ hover $\rightarrow$ loaded $\rightarrow$ idle on reset) work for all supported file types. \hfill\textit{(FL-010)}
  \item \textbf{RT-005}: Voxel Comparison info label updates correctly after every augmentation. \hfill\textit{(RC-009)}
\end{enumerate}

% =============================================================================
\section{Test Execution Summary}
% =============================================================================

\begin{longtable}{p{1.2cm}p{5.5cm}p{1.5cm}p{1.5cm}p{1.5cm}p{2cm}}
\toprule
\textbf{ID} & \textbf{Title} & \textbf{Priority} & \textbf{Pass/Fail} & \textbf{Tester} & \textbf{Date} \\
\midrule
FL-001 & Drag \& drop .ply & High & & & \\
FL-002 & Drag \& drop .obj & High & & & \\
FL-003 & Drag \& drop .npy & High & & & \\
FL-004 & Drag \& drop unsupported & Medium & & & \\
FL-005 & Browse dialog filters & High & & & \\
FL-006 & Browse select .ply & High & & & \\
FL-007 & Browse select .npy & High & & & \\
FL-008 & Click Drop Zone opens dialog & Medium & & & \\
FL-009 & Load Mesh button re-loads & Low & & & \\
FL-010 & Drag visual state transitions & Medium & & & \\
FL-011 & Load corrupted file & High & & & \\
\addlinespace
FI-001 & Info panel mesh stats & Medium & & & \\
FI-002 & Info panel voxel stats & Medium & & & \\
FI-003 & Info panel after voxelization & Medium & & & \\
\addlinespace
PL-001 & Mesh $\rightarrow$ Voxels (default) & High & & & \\
PL-002 & Mesh $\rightarrow$ Voxels (custom) & Medium & & & \\
PL-003 & Mesh $\rightarrow$ Voxels (boundary) & Low & & & \\
PL-004 & Dilate (default) & High & & & \\
PL-005 & Dilate (various iterations) & Medium & & & \\
PL-006 & Voxels $\rightarrow$ Mesh & High & & & \\
PL-007 & Full pipeline (end-to-end) & High & & & \\
PL-008 & Pipeline button states & High & & & \\
\addlinespace
AU-001 & Deformation each axis & High & & & \\
AU-002 & Deformation min/max factor & Medium & & & \\
AU-003 & Erosion each axis & High & & & \\
AU-004 & Erosion min/max increment & Medium & & & \\
AU-005 & Rotation zero angles & Low & & & \\
AU-006 & Rotation non-zero angles & Medium & & & \\
AU-007 & Rotation max angles & Low & & & \\
AU-008 & Fracture return\_both=False & High & & & \\
AU-009 & Fracture return\_both=True & Medium & & & \\
AU-010 & Fracture max\_pos boundary & Low & & & \\
AU-011 & Save Deformed Voxels & High & & & \\
AU-012 & Augmentation chain & Medium & & & \\
AU-013 & Buttons disabled without voxels & Medium & & & \\
\addlinespace
VS-001 & View Mesh & High & & & \\
VS-002 & View Voxels & High & & & \\
VS-003 & View empty voxels & Low & & & \\
VS-004 & Color picker & Medium & & & \\
VS-005 & Save Mesh .ply & High & & & \\
VS-006 & Save Mesh .obj & Medium & & & \\
VS-007 & Save Voxel .npy & High & & & \\
VS-008 & Default save filenames & Low & & & \\
\addlinespace
RC-001 & Reconstruct from Voxels & High & & & \\
RC-002 & Reconstruct from .npy & High & & & \\
RC-003 & Compare Both (meshes) & High & & & \\
RC-004 & Compare Original Only & Medium & & & \\
RC-005 & Compare Reconstructed Only & Medium & & & \\
RC-006 & Compare Both (voxels) & High & & & \\
RC-007 & Voxel comparison Original Only & Medium & & & \\
RC-008 & Voxel comparison Current Only & Medium & & & \\
RC-009 & Voxel comparison info label & Medium & & & \\
RC-010 & Save Reconstructed Mesh & High & & & \\
\addlinespace
TS-001 & t-SNE default params & High & & & \\
TS-002 & t-SNE custom params & Medium & & & \\
TS-003 & t-SNE boundary values & Low & & & \\
TS-004 & t-SNE empty voxels error & Medium & & & \\
TS-005 & Save t-SNE image & Medium & & & \\
TS-006 & Close t-SNE dialog & Low & & & \\
TS-007 & t-SNE button state & Medium & & & \\
\addlinespace
D2-001 & 2D single XY view & High & & & \\
D2-002 & 2D single XZ view & Medium & & & \\
D2-003 & 2D single YZ view & Medium & & & \\
D2-004 & 2D different colors & Low & & & \\
D2-005 & 2D different markers & Low & & & \\
D2-006 & 2D marker size & Low & & & \\
D2-007 & 2D comparison & High & & & \\
D2-008 & 2D comparison axes & Medium & & & \\
D2-009 & Save 2D image & Medium & & & \\
D2-010 & 2D button states & Medium & & & \\
\addlinespace
RA-001 & Reset All (full state) & High & & & \\
RA-002 & Reset All (npy only) & Medium & & & \\
RA-003 & About dialog display & Low & & & \\
RA-004 & About dialog close & Low & & & \\
\addlinespace
SE-001 & Initial status bar message & Low & & & \\
SE-002 & Status bar updates & Medium & & & \\
SE-003 & Error dialog on failure & Medium & & & \\
SE-004 & Window title and icon & Low & & & \\
\bottomrule
\end{longtable}

\vspace{1cm}
\begin{center}
\large\textbf{Total Test Cases: 67}
\end{center}

\vspace{0.5cm}
\begin{tabular}{lrr}
\toprule
\textbf{Priority} & \textbf{Count} \\
\midrule
High   & 31 \\
Medium & 27 \\
Low    & 9  \\
\bottomrule
\end{tabular}

\end{document}
